\subsection{Circumstances Often Segment Better Than Identity}

The exact same human being exhibits completely different purchasing behavior depending upon the acute circumstance. A developer casually browsing tutorials on a Sunday afternoon exhibits high price sensitivity and low urgency. The same developer facing an executive audit of model reliability on Thursday morning suddenly exhibits urgent, high-stakes demand. A parent buying food on a frantic weekday morning prioritizes speed and convenience; on a weekend holiday, they prioritize health reassurance and culinary craft.

Segmenting by \textbf{circumstance and buying trigger} reveals concentrated demand that demographic averages conceal. However, not every situational segment is economically viable. A viable segment requires an identifiable trigger, sufficient transaction volume, distinct purchasing behavior, and a practical, affordable route to reach the buyer at the moment the trigger fires.

\subsection{The Switching Story vs. The Wish List}

Customer wish lists are lists of free aspirations. A switching story is an forensic reconstruction of a costly behavioral change:
\begin{center}
\small
Initial discomfort $\rightarrow$ Triggering incident $\rightarrow$ Passive search $\rightarrow$ Active comparison $\rightarrow$ Shortlist $\rightarrow$ Anxieties $\rightarrow$ Decisive catalyst $\rightarrow$ Transaction $\rightarrow$ Onboarding $\rightarrow$ Usage $\rightarrow$ Outcome.
\end{center}

In the switching timeline, search for \textbf{struggle moments}: ``I attempted 3 tutorials over 2 weekends, but our retrieval pipeline failed completely in the release demo.'' ``Our dean mandated an AI elective, but the faculty workload was already at 100\% capacity.'' ``We evaluated 2 software vendors, but corporate legal rejected their data retention clauses.''

These struggle moments reveal the real operational bottlenecks. The objective is not coaxing drama from routine decisions, but uncovering the causal friction that forces an actor to abandon a comfortable default.

\subsection{Nonconsumption as the Dominant Competitor}

A customer who never buys has not necessarily chosen a branded competitor. They may be tolerating the problem, lacking money, unaware that an alternative is possible, or blocked by corporate governance. \textbf{Nonconsumption is a collection of distinct causal mechanisms, not a single monolithic market.}

Distinguish constrained nonconsumers: actors who actively desire a solution but are blocked by price, physical access, or regulatory permission: from unconcerned nonconsumers, who simply do not value the job enough to act. The former can be unlocked through lower-friction delivery, modular pilots, or employer-funded sponsorships. The latter do not represent a commercial market today.

Interview people who investigated a category and explicitly chose to do nothing. Their narratives dismantle the common founder fantasy that every person with a problem is an eager customer in waiting.

\subsection{Desired Outcome vs. Proposed Feature}

Suppose 5 learners request recorded lecture videos. One needs temporal flexibility around family duties. Another fears missing dense mathematical derivations in real time. A third wants to review code syntax repeatedly. A fourth suffers from erratic internet connectivity. A fifth feels social anxiety when questioned live.

A founder who notes ``5 votes for recordings'' builds an undifferentiated video repository that fails all 5. Some needed downloadable offline materials; some needed interactive self-paced debugging harnesses; some needed psychological safety in live seminars.

Translate customer feature requests into \textbf{desired outcome statements}: what specific capability breaks without the feature, and what alternative intervention would achieve the outcome with less operational overhead?

\subsection{The Boundaries of Jobs Theory}

Jobs-to-be-done is an indispensable lens for understanding behavioral motivation, but it does not replace the laws of industrial organization, balance-sheet economics, or institutional power.

Two individuals can experience the identical job: ``learn production AI systems'': while 1 possesses ₹50,000 in disposable savings and the other is deeply in debt. Three engineers may struggle with identical deployment bottlenecks, yet only their division director holds the corporate credit card. An administrator may recognize an urgent educational job for students, yet state university regulations prohibit hiring external teaching contractors.

Never reduce strategic analysis to an emotionally moving jobs narrative. A complete customer model combines the job with \textbf{budget ownership, legal procurement, competitive alternatives, cash flow cycles, price sensitivity, and structural market power}.
